% Gemini theme
% https://github.com/anishathalye/gemini
%
% NeurIPS 2026 poster: Mixture-of-Control (MoC).
% Main-conference format: 140 cm x 90 cm, landscape.

\documentclass[final]{beamer}

% ====================
% Packages
% ====================

\usepackage[size=custom,width=140,height=90,scale=1.0]{beamerposter}
\usetheme{gemini}
\usecolortheme{neurips}
\usepackage{graphicx}
\usepackage{booktabs}
\usepackage{tabularx}
\usepackage{array}
\usepackage{colortbl}
\usepackage{amsmath}
% Sans-serif math matched to the Lato body text, so formulas, tables and
% prose share one typeface family.
\usepackage{unicode-math}
\setmathfont{LeteSansMath.otf}[Scale=MatchLowercase]
% Upright math text (\mathbf, \text-like operators) also in Lato.
\setmathrm{Lato}
\usepackage{algpseudocode}
\usepackage{tcolorbox}
\tcbuselibrary{skins}
\usepackage{tikz}
\usetikzlibrary{positioning}
\usepackage{qrcode}
\usepackage{anyfontsize}

% ====================
% Paper metadata
% ====================

\newcommand{\paperid}{NeurIPS 2026 \#148178}
\newcommand{\paperurl}{https://neurips.cc/virtual/2026/loc/sydney/poster/148178}
\newcommand{\arxivurl}{https://arxiv.org/abs/2606.31397}
\newcommand{\figdir}{MoC-project-page/static/images}

\title{Mixture-of-Control: State-Aware Fine-Tuning for Transformer-based Models}
\author[Nguyen, Luu, Pham, Tran]{\texorpdfstring{%
  Duc-Anh Nguyen\inst{1,2}\textsuperscript{*,\textdagger}
  \and Tien Ngoc Luu\inst{1}\textsuperscript{*}
  \and Tung Pham\inst{1}
  \and Toan Tran\inst{1}%
}{Duc-Anh Nguyen, Tien Ngoc Luu, Tung Pham, Toan Tran}}
\institute[Qualcomm AI Research; University of Virginia]{%
  \inst{1} Qualcomm AI Research\textsuperscript{\textdaggerdbl}
  \samelineand \inst{2} University of Virginia
  \\[-0.1ex]\footnotesize
  * Equal contribution \quad
  \textdagger\ Work done during an internship at Qualcomm \quad
  Corresponding author: Toan Tran \quad
  \textdaggerdbl\ Qualcomm AI Research is an initiative of Qualcomm Technologies, Inc.}

% ====================
% Layout
% ====================

% Four separators and three columns fill the paper width:
% 4*\sepwidth + 2*\colwidth + \widewidth = \paperwidth
\newlength{\sepwidth}
\newlength{\colwidth}
\newlength{\widewidth}
\setlength{\sepwidth}{0.018\paperwidth}
\setlength{\colwidth}{0.248\paperwidth}
\setlength{\widewidth}{0.432\paperwidth}
\newcommand{\separatorcolumn}{\begin{column}{\sepwidth}\end{column}}

% ====================
% Colors and reusable visual elements
% ====================

\definecolor{mocup}{HTML}{1B7F4B}
\colorlet{mocrow}{neuripslavender}
\colorlet{grouprow}{neuripsrule!55}
\colorlet{owncolor}{neuripspurple}
\definecolor{routecolor}{HTML}{185FA5}
\definecolor{pagebg}{HTML}{F4F3F8}

\newcommand{\own}[1]{{\color{owncolor}#1}}
\newcommand{\route}[1]{{\color{routecolor}#1}}

\newcommand{\pill}[2]{%
  \tikz[baseline=(label.base)]\node[
    rounded corners=1.2ex,
    fill=#1,
    text=white,
    inner xsep=1.1ex,
    inner ysep=0.45ex,
    font=\bfseries\small
  ] (label) {#2};%
}

% Numbered "gap" card used in the motivation block.
\newcommand{\gapcard}[3]{%
  \begin{tikzpicture}
    \node[fill=neuripsrose,rounded corners=0.6cm,inner sep=0.8cm,
          text width=\linewidth-1.6cm,align=left] (c) {%
      {\color{neuripscoral}\bfseries\Large #1}\quad{\bfseries #2}\\[0.4ex]
      \small #3};
  \end{tikzpicture}%
}

% Rounded callout for a single statement.
\newcommand{\callout}[2][neuripslavender]{%
  \begin{tikzpicture}
    \node[fill=#1,rounded corners=0.6cm,inner sep=0.7cm,
          text width=\linewidth-1.4cm,align=center] {\Centering #2};
  \end{tikzpicture}%
}

% Boxed theorem statement for the theory block.
\newtcolorbox{thmbox}[2]{%
  enhanced,
  colback=white, colframe=neuripspurple!75,
  colbacktitle=neuripslavender, coltitle=neuripsnavy,
  boxrule=1.6pt, arc=3mm,
  left=7mm, right=7mm, top=3mm, bottom=3mm,
  toptitle=1.5mm, bottomtitle=1.5mm,
  fonttitle=\Raleway\bfseries,
  title={{\color{neuripsmagenta}#1}\enspace#2},
  before skip=2.2ex, after skip=1.2ex}

% Algorithm box with a dark title bar.
\newtcolorbox{algbox}[1]{%
  enhanced,
  colback=neuripslavender!40, colframe=neuripsnavy,
  colbacktitle=neuripsnavy, coltitle=white,
  boxrule=1.6pt, arc=3mm,
  left=8mm, right=8mm, top=4mm, bottom=4mm,
  toptitle=2mm, bottomtitle=2mm,
  fonttitle=\Raleway\bfseries,
  title={#1}}

% Muted, right-aligned algorithm comments.
\algrenewcommand{\algorithmiccomment}[1]{\hfill{\color{neuripsmuted}\small$\triangleright$~#1}}
\algrenewcommand{\algorithmicrequire}{\textbf{Input:}}
\algrenewcommand{\algorithmicensure}{\textbf{Output:}}

% Tighter display math; size commands reset these skips, so hook into them.
\newcommand{\tightdisplay}{%
  \setlength{\abovedisplayskip}{0.5ex}\setlength{\belowdisplayskip}{0.5ex}%
  \setlength{\abovedisplayshortskip}{0.3ex}\setlength{\belowdisplayshortskip}{0.3ex}}
\makeatletter
\g@addto@macro\normalsize{\tightdisplay}
\g@addto@macro\small{\tightdisplay}
\makeatother

% Centered equation line without display-math skips.
\newcommand{\eqline}[1]{\par{\centering$\displaystyle #1$\par}}

% Small delta-vs-LoRA tag inside the main table.
\newcommand{\up}[1]{\,{\color{mocup}\scriptsize$\uparrow$#1}}

% ====================
% Header and footer
% ====================

% ---- White headline: dark title text, official logos placed directly.
\setbeamercolor{headline}{fg=neuripsnavy,bg=white}
\setbeamercolor{headline title}{fg=neuripsnavy}
\setbeamercolor{headline author}{fg=neuripsink}
\setbeamercolor{headline institute}{fg=neuripsmuted}
\setbeamercolor{headline rule}{bg=neuripspurple}
\setbeamercolor{background canvas}{bg=pagebg}

\logoleft{%
  \tikz[baseline=(q.south)]{%
    \node (q) {\includegraphics[height=2.2cm]{logos/qualcomm.pdf}};
    \node[below=1.2cm of q] {\includegraphics[height=3.2cm]{logos/uva_wordmark.pdf}};}}
\logoright{\includegraphics[height=6.4cm]{logos/neurips.pdf}}

% ---- Card-style blocks: white rounded card with a colored title band.
\tcbset{
  poster card/.style 2 args={
    enhanced, arc=5mm,
    colback=#2, colframe=#1, boxrule=1.2pt,
    colbacktitle=#1, coltitle=white,
    fonttitle={\Raleway\large\bfseries},
    toptitle=3.5mm, bottomtitle=3.5mm, lefttitle=9mm,
    left=9mm, right=9mm, top=5mm, bottom=6mm,
    before skip=0pt, after skip=2.4ex}}
\setbeamertemplate{block begin}{%
  \begin{tcolorbox}[poster card={neuripspurple}{white},title={\insertblocktitle}]%
  \usebeamerfont{block body}\justifying\setlength{\parskip}{1ex}}
\setbeamertemplate{block end}{\end{tcolorbox}}
\setbeamertemplate{block alerted begin}{%
  \begin{tcolorbox}[poster card={neuripsmagenta}{neuripsrose},title={\insertblocktitle}]%
  \usebeamerfont{block body}\justifying\setlength{\parskip}{1ex}}
\setbeamertemplate{block alerted end}{\end{tcolorbox}}

\footercontent{%
  \paperid \hfill
  \href{\arxivurl}{arXiv:2606.31397} \hfill
  \href{\paperurl}{neurips.cc/virtual/2026/poster/148178}%
}

% ====================
% Body
% ====================

\begin{document}

\begin{frame}[t]

\vspace{1ex}

\begin{columns}[t]
\separatorcolumn

% ====================
% Column 1: motivation, problem, theory
% ====================

\begin{column}{\colwidth}

  \begin{block}{1. Motivation: two gaps in efficient fine-tuning}

    \pill{neuripspurple}{BACKGROUND}\;
    \textbf{State-based} fine-tuning updates lightweight controls on the
    \emph{states}, not the weights: large memory savings with high parameter
    efficiency.

    \vspace{1.2ex}
    \gapcard{1}{Control stays local to each block}{LoRA variants and
      state-based methods such as Parallel Control adapt each block in
      isolation $\Rightarrow$ weak coordination across depth.}

    \vspace{0.9ex}
    \gapcard{2}{Cross-layer communication is costly}{MoLEx mixes layers by
      \textbf{re-running pretrained blocks}. On Qwen2.5-14B:
      \textbf{124.8\,GB / 20.3\,h} vs.\ 62.1\,GB / 9.4\,h for LoRA.}

  \end{block}

  \begin{block}{2. Problem definition}

    A transformer forward pass is a \textbf{discretized feedback system}:
    frozen block $f_t$ = intrinsic dynamics, low-rank adapter $g_t$ = control.

    \heading{Parallel Control (block-local)}
    \eqline{x_{t+1}=x_t+f_t(W_tx_t)+g_t(x_t)}

    \vspace{0.8ex}
    \heading{Generalized state-based update}
    \eqline{x_{t+1}=x_t+f_t(W_tx_t)+\alpha\,\own{g_t(x_t)}+(1-\alpha)\,\route{g_{h_t}(x_t)}}
    \small $h_t=t$ recovers local tuning; $h_t\neq t$ communicates across
    blocks through \textbf{low-rank controls}, not full-block reuse.

    \vspace{1.2ex}
    \callout{\itshape How can state-based fine-tuning capture deep cross-layer
      interactions \textbf{without sacrificing memory efficiency}?}

  \end{block}

  \begin{block}{3. Theory: why routing control helps}

    Under mild assumptions, MoC \textbf{keeps MoLEx's expressivity} and is
    provably better behaved than block-local Parallel Control.

    \begin{thmbox}{Theorem 4.1}{(Expressivity via state controls)}
      If $\alpha>\tfrac12$ and $z_1=x_1$, there exist controls $g_k$ that
      reproduce the MoLEx trajectory: exactly for linear $f_k$, and with
      vanishing error for smooth nonlinear $f_k$ under small perturbations
      $\|\Delta W'_t\|<\xi$:
      \eqline{z_t=x_t\ \ \forall t\ \text{(linear)},\qquad
        \mathbb{E}\|z_t-x_t\|^2\xrightarrow{\ \xi\to0\ }0}
    \end{thmbox}

    \begin{thmbox}{Proposition}{(Stability)}
      With routed magnitude $\beta_K(x_t)=\sum_{k\in S_K}\pi_k\beta_k$ and
      $\mathbb{E}[\beta_K\mid x_t]<\beta_t$ (encouraged by load balancing),
      MoC's expected one-step growth is strictly smaller:
      \eqline{1+L_f+\big[\alpha\beta_t+(1-\alpha)\,\mathbb{E}[\beta_K]\big]\gamma
        \;<\;1+L_f+\beta_t\gamma}
    \end{thmbox}

    \begin{thmbox}{Corollary}{(Gradient flow)}
      Each layerwise Jacobian bound tightens, $M_t\le M_t^{c}$, so the
      backpropagated product is better conditioned:
      \eqline{\prod_{t=1}^{L}\|J_t\|\;\le\;\prod_{t=1}^{L}M_t\;\le\;\prod_{t=1}^{L}M_t^{c}}
    \end{thmbox}

    \begin{thmbox}{Proposition}{(Representation)}
      Mixing routed experts brings the control closer to the task-optimal
      signal $\Delta W^{*}x$ than block-local control alone:
      \eqline{\bar\epsilon^{2}\;\le\;\alpha\,\bar\epsilon_c^{2}+(1-\alpha)\,\bar\epsilon_K^{2}\;<\;\bar\epsilon_c^{2}}
    \end{thmbox}

  \end{block}

\end{column}

\separatorcolumn

% ====================
% Column 2: method, algorithm, conclusion
% ====================

\begin{column}{\colwidth}

  \begin{block}{4. Method: control states as experts}

    \begin{figure}
      \centering
      % Trim the PNG's uneven side margins (ink spans x=59..2283 of 2400 px)
      % so the drawing is optically centered in the column.
      \includegraphics[width=\linewidth,trim=40 0 100 0,clip]{\figdir/MoC_architecture.png}
      \caption{At block $i$, frozen $f_i$ and local control $g_i$ run in
        parallel. A \textbf{shared gate} picks Top-$K$ controls from
        \emph{all} layers and fuses them with the local one ($K{=}1$ shown).}
    \end{figure}

    \vspace{1ex}
    \heading{MoC update}
    \eqline{x_{t+1}=x_t+f_t(W_tx_t)+\alpha\,\own{g_t(x_t)} +(1-\alpha)\sum_{k=1}^{L}\pi_k(x_t)\,\route{g_k(x_t)}}
    \vspace{0.8ex}
    \heading{Sparse routing with a shared gate}
    \eqline{\pi_k(x_t)=\operatorname{softmax}\!\big(\operatorname{TopK}(\mathbf{W}_{\text{gate}}\,x_t)\big)_k}

    \vspace{1.2ex}
    \begin{columns}[t,totalwidth=\linewidth]
      \begin{column}{0.49\linewidth}
        \begin{itemize}
          \item \textbf{Input:} hidden state $x_t$ at block $t$.
          \item \textbf{Expert pool:} the $L$ controls $\{g_k\}$ of all blocks.
          \item \textbf{Gate:} one $\mathbf{W}_{\text{gate}}$ shared across layers.
        \end{itemize}
      \end{column}
      \begin{column}{0.49\linewidth}
        \begin{itemize}
          \item \textbf{Objective:} $\mathcal{L}=\mathcal{L}_{\text{task}}+\lambda\,\mathcal{L}_{\text{balance}}$.
          \item \textbf{Training:} batch-level routing (mean / mode of token scores).
          \item \textbf{Inference:} each sample routed from its own prefix.
        \end{itemize}
      \end{column}
    \end{columns}

  \end{block}

  \begin{block}{5. Algorithm: MoC forward pass}

    \begin{algbox}{Algorithm 1\enspace MoC forward pass at block $t$}
      \setlength{\baselineskip}{1.45\baselineskip}
      \begin{algorithmic}[1]
        \Require state $x_t$; frozen block $f_t$; controls $\{g_k\}_{k=1}^{L}$; gate $\mathbf{W}_{\text{gate}}$; $K$, $\alpha$
        \Ensure next state $x_{t+1}$
        \State $p_{\text{own}}\gets \own{g_t(x_t)}$ \Comment{layer-dedicated expert}
        \State $\pi\gets\operatorname{softmax}(\mathbf{W}_{\text{gate}}\,x_t)$ \Comment{dense routing}
        \State $S_K\gets\operatorname{TopK}(\pi,K)$ \Comment{expert selection}
        \State $\pi_i\gets\pi_i/\sum_{j\in S_K}\pi_j,\ \ i\in S_K$ \Comment{re-normalize}
        \State $p_{\text{route}}\gets\sum_{i\in S_K}\pi_i\,\route{g_i(x_t)}$ \Comment{sparse aggregation}
        \State $g_{\text{moc}}\gets\alpha\,p_{\text{own}}+(1-\alpha)\,p_{\text{route}}$ \Comment{control signal}
        \State \Return $x_{t+1}\gets x_t+f_t(x_t)+g_{\text{moc}}$ \Comment{$f_t$ runs once}
      \end{algorithmic}
    \end{algbox}

  \end{block}

  \begin{alertblock}{Key findings \& conclusion}
    \begin{columns}[T,totalwidth=\linewidth]
      \begin{column}{0.74\linewidth}
        \begin{itemize}
          \item \textbf{Novelty:} block-wise control states become a
            cross-layer \textbf{sparse MoE pool} with one shared gate.
          \item \textbf{Accuracy:} best average on all 4 LLMs, GLUE, and
            VTAB-1K; up to \textbf{+5.9} over LoRA.
          \item \textbf{Efficiency:} pretrained blocks run once
            $\Rightarrow$ \textbf{$\sim$half the memory} of MoLEx, throughput
            near Control.
          \item \textbf{Theory:} keeps MoLEx expressivity with better
            stability, gradients, and representation error.
          \item \textbf{Limitation:} like other state-based methods, MoC
            is still slower at inference than fully merged weight adapters.
        \end{itemize}
      \end{column}
      \begin{column}{0.24\linewidth}
        \centering
        \qrcode[height=6cm]{\arxivurl}\\[0.4ex]
        \footnotesize\bfseries arXiv:2606.31397
      \end{column}
    \end{columns}
  \end{alertblock}

\end{column}

\separatorcolumn

% ====================
% Column 3: results, ablations, throughput
% ====================

\begin{column}{\widewidth}

  \begin{block}{6. Main results: commonsense reasoning on four LLM backbones}

    \begin{table}
      \centering
      \footnotesize
      \setlength{\tabcolsep}{0.32em}
      \renewcommand{\arraystretch}{1.15}
      \begin{columns}[t,totalwidth=\linewidth]
        \begin{column}{0.495\linewidth}
          \centering
          \begin{tabular}{@{}l c cccccccc >{\columncolor{neuripsrule!30}}c cc@{}}
            \toprule
            \textbf{Method} & \textbf{\%Par.} & \textbf{BoolQ} & \textbf{PIQA} & \textbf{SIQA} & \textbf{HellaS.} & \textbf{WinoG.} & \textbf{ARC-e} & \textbf{ARC-c} & \textbf{OBQA} & \textbf{Avg$\uparrow$} & \textbf{h$\downarrow$} & \textbf{GB$\downarrow$} \\
            \midrule
            \rowcolor{grouprow}\multicolumn{13}{l}{\textbf{LLaMA2-7B}} \\
            LoRA$^{*}$ & 0.83 & 69.8 & 79.9 & 79.5 & 83.6 & 82.6 & 79.8 & 64.7 & 81.0 & 77.6 & 6.1 & 48.0 \\
            DoRA$^{*}$ & 0.85 & 71.8 & \textbf{83.7} & 76.0 & 89.1 & 82.6 & \textbf{83.7} & 68.2 & \textbf{82.4} & 79.7 & 10.5 & 80.2 \\
            MoLEx & 0.85 & 69.7 & 81.9 & 78.4 & 91.3 & 81.5 & 81.5 & 65.8 & 80.8 & 78.9 & 9.8 & 77.4 \\
            FacT\_TK & 0.51 & 69.8 & 80.6 & 77.8 & 90.8 & 81.8 & 79.7 & 65.4 & 78.6 & 78.1 & 6.0 & 48.5 \\
            StructLoRA & 0.83 & 70.9 & 82.0 & 78.6 & 91.6 & 82.1 & 81.6 & 66.6 & 80.6 & 79.3 & 8.5 & 71.6 \\
            FlexLoRA & 0.83 & 69.5 & 81.2 & 78.1 & 92.9 & 78.7 & 82.8 & 66.5 & 77.0 & 78.3 & 7.0 & 54.8 \\
            Control$^{*}$ & 0.50 & \textbf{72.3} & 82.5 & 79.2 & 89.1 & 83.1 & 83.0 & 68.5 & 79.0 & 79.6 & 4.2 & 42.6 \\
            \rowcolor{mocrow}\textbf{MoC} & 0.51 & 71.8 & 83.2 & \textbf{79.6} & \textbf{93.3} & \textbf{84.2} & 83.6 & \textbf{68.6} & 80.4 & \textbf{80.5}\up{2.9} & 5.5 & 47.6 \\
            \midrule
            \rowcolor{grouprow}\multicolumn{13}{l}{\textbf{LLaMA3-8B}} \\
            LoRA$^{*}$ & 0.71 & 70.8 & 85.2 & 79.9 & 91.7 & 84.3 & 84.2 & 71.2 & 79.0 & 80.8 & 7.5 & 59.1 \\
            DoRA$^{*}$ & 0.71 & 74.6 & 89.3 & 79.9 & 95.5 & 85.6 & 90.5 & 80.4 & 85.8 & 85.2 & 10.5 & 69.1 \\
            MoLEx & 0.72 & 70.1 & 86.2 & 80.4 & 93.2 & 84.1 & 84.8 & 71.3 & 82.2 & 81.5 & 11.0 & 94.9 \\
            FacT\_TK & 0.43 & 70.7 & 86.0 & 79.5 & 94.2 & 84.2 & 86.5 & 73.3 & 81.8 & 82.0 & 7.7 & 57.6 \\
            StructLoRA & 0.71 & 70.9 & 86.6 & 79.5 & 93.9 & 84.0 & 87.0 & 73.9 & 82.6 & 82.3 & 9.8 & 75.4 \\
            FlexLoRA & 0.71 & 73.4 & 88.5 & 78.6 & 95.3 & 84.4 & 90.2 & 79.2 & 85.6 & 84.4 & 8.6 & 62.4 \\
            Control$^{*}$ & 0.42 & 74.1 & 87.8 & 80.7 & 95.5 & 86.0 & 90.8 & 80.0 & 87.8 & 85.3 & 4.4 & 54.0 \\
            \rowcolor{mocrow}\textbf{MoC} & 0.43 & \textbf{75.5} & \textbf{89.5} & \textbf{82.0} & \textbf{96.6} & \textbf{88.8} & \textbf{91.5} & \textbf{81.6} & \textbf{88.0} & \textbf{86.7}\up{5.9} & 6.5 & 56.7 \\
            \bottomrule
          \end{tabular}
        \end{column}
        \begin{column}{0.495\linewidth}
          \centering
          \begin{tabular}{@{}l c cccccccc >{\columncolor{neuripsrule!30}}c cc@{}}
            \toprule
            \textbf{Method} & \textbf{\%Par.} & \textbf{BoolQ} & \textbf{PIQA} & \textbf{SIQA} & \textbf{HellaS.} & \textbf{WinoG.} & \textbf{ARC-e} & \textbf{ARC-c} & \textbf{OBQA} & \textbf{Avg$\uparrow$} & \textbf{h$\downarrow$} & \textbf{GB$\downarrow$} \\
            \midrule
            \rowcolor{grouprow}\multicolumn{13}{l}{\textbf{Mistralv3-7B}} \\
            LoRA & 0.78 & 68.7 & 83.5 & 77.2 & 90.3 & 82.5 & 85.7 & 73.4 & 80.5 & 80.2 & 6.5 & 58.4 \\
            DoRA & 0.79 & 70.3 & 84.5 & 77.8 & 89.9 & 83.1 & 87.1 & 72.2 & 82.4 & 80.9 & 10.1 & 70.4 \\
            MoLEx & 0.80 & 71.4 & 85.4 & 78.2 & 90.3 & 83.5 & 84.1 & 75.0 & 82.7 & 81.3 & 10.5 & 82.3 \\
            FacT\_TK & 0.48 & 70.5 & 86.4 & 80.3 & 94.1 & 84.2 & 83.9 & 71.0 & 84.4 & 81.9 & 6.3 & 56.0 \\
            StructLoRA & 0.78 & 70.0 & 83.2 & 77.4 & 89.9 & 81.7 & 87.2 & 74.8 & 81.4 & 80.7 & 9.2 & 70.1 \\
            FlexLoRA & 0.78 & 73.9 & 88.4 & 79.2 & \textbf{95.9} & 86.0 & \textbf{90.2} & 78.8 & 86.8 & 84.9 & 7.5 & 61.3 \\
            Control & 0.46 & 70.6 & 86.9 & 79.0 & 94.2 & 84.8 & 87.9 & 76.5 & 84.8 & 83.1 & 4.3 & 53.2 \\
            \rowcolor{mocrow}\textbf{MoC} & 0.48 & \textbf{74.2} & \textbf{88.9} & \textbf{81.0} & 95.6 & \textbf{87.9} & 88.9 & \textbf{79.5} & \textbf{87.8} & \textbf{85.5}\up{5.3} & 5.5 & 55.6 \\
            \midrule
            \rowcolor{grouprow}\multicolumn{13}{l}{\textbf{Qwen2.5-14B}} \\
            LoRA & 0.63 & 73.4 & 90.4 & 82.2 & 96.0 & 88.6 & 94.7 & 86.1 & 90.6 & 87.7 & 9.4 & 62.1 \\
            DoRA & 0.64 & 75.1 & 90.1 & 82.0 & 96.3 & 89.4 & 95.0 & 88.2 & 93.0 & 88.6 & 14.3 & 94.7 \\
            MoLEx & 0.64 & 73.5 & 90.5 & 81.9 & 95.5 & 88.2 & 94.4 & 85.8 & 92.2 & 87.8 & 20.3 & 124.8 \\
            FacT\_TK & 0.43 & 74.7 & 90.1 & 82.3 & 96.2 & 89.6 & 94.7 & 88.1 & 92.4 & 88.5 & 8.0 & 61.5 \\
            StructLoRA & 0.63 & 74.8 & 89.5 & 81.0 & 95.8 & 89.3 & 94.7 & 87.4 & 92.4 & 88.1 & 13.5 & 77.3 \\
            FlexLoRA & 0.63 & 74.2 & 91.4 & 82.0 & 96.4 & 87.7 & \textbf{96.7} & 90.0 & 95.4 & 89.2 & 10.5 & 67.4 \\
            Control & 0.43 & 72.8 & 88.8 & 81.3 & 94.6 & 87.5 & 94.4 & 86.9 & 91.2 & 87.2 & 6.5 & 58.7 \\
            \rowcolor{mocrow}\textbf{MoC} & 0.43 & \textbf{75.5} & \textbf{91.8} & \textbf{82.4} & \textbf{96.8} & \textbf{91.0} & 96.6 & \textbf{90.2} & \textbf{96.4} & \textbf{90.1}\up{2.4} & 8.0 & 61.3 \\
            \bottomrule
          \end{tabular}
        \end{column}
      \end{columns}
      \caption{Accuracy (\%) on 8 commonsense datasets. \textbf{Bold}: best per
        column within each backbone; {\color{mocup}$\uparrow$}: average gain over
        LoRA; $^{*}$: previously published numbers. \%Par.: trainable parameters;
        h: training hours; GB: peak training GPU memory. MoC also achieves the
        best average on GLUE (RoBERTa, DeBERTa) and VTAB-1K (ViT-Tiny, ViT-Base).}
    \end{table}

  \end{block}

  \begin{block}{7. Ablations: what makes MoC work}

    \begin{figure}
      \centering
      \includegraphics[width=0.88\linewidth]{\figdir/ablation_hyperparams.png}
      \caption{(a) Cost grows moderately with \#experts (Mistralv3-7B).
        (b) Best: local + 1 routed expert. (c) Robust to $\alpha$, peak at
        0.95. (d) Load balancing peaks at $\lambda{=}0.01$.}
    \end{figure}

    \begin{columns}[t,totalwidth=\linewidth]
      \begin{column}{0.53\linewidth}
        \centering\small\renewcommand{\arraystretch}{1.2}
        \begin{tabular}{l c c}
          \toprule
          \textbf{Communication} & \textbf{Acc.} & \textbf{GB} \\
          \midrule
          Control (local) & 83.1 & 53.2 \\
          Uniform Control & 84.3 & 53.8 \\
          Fixed Control & 84.2 & 55.0 \\
          MoC w/o load bal. & 84.4 & 55.4 \\
          \rowcolor{mocrow}\textbf{MoC} & \textbf{85.5} & 55.6 \\
          \bottomrule
        \end{tabular}\\[0.3ex]
        {\color{neuripsmuted}Mistralv3-7B}
      \end{column}
      \begin{column}{0.45\linewidth}
        \centering\small\renewcommand{\arraystretch}{1.2}
        \begin{tabular}{l c c}
          \toprule
          \textbf{Variant} & \textbf{L3} & \textbf{M3} \\
          \midrule
          ATTN & 85.9 & 84.9 \\
          FFN & 86.2 & 85.1 \\
          A+F w/o share & 86.4 & 84.9 \\
          \rowcolor{mocrow}\textbf{A+F shared} & \textbf{86.7} & \textbf{85.5} \\
          \bottomrule
        \end{tabular}\\[0.3ex]
        {\color{neuripsmuted}L3: LLaMA3-8B, M3: Mistralv3-7B}
      \end{column}
    \end{columns}


    \vspace{0.6ex}
    {\small Every communication variant beats local Control; adaptive
    routing + load balancing adds \textbf{+1.1}, and sharing one pool across
    ATTN and FFN controls works best.}

  \end{block}

  \begin{block}{8. Inference throughput}

    \begin{figure}
      \centering
      \includegraphics[width=0.80\linewidth]{\figdir/inference_throughput_moc_comparison.png}
      \caption{Tokens/s vs.\ batch size. Shaded: gap to fully merged
        weight adapters (upper bound). MoC tracks Control and stays well
        above MoLEx, which re-runs pretrained blocks.}
    \end{figure}

  \end{block}

\end{column}

\separatorcolumn
\end{columns}
\end{frame}

\end{document}
